\section{Тензорные произведения}

Пусть $V_1$, \dots, $V_m$ --- (конечномерные) векторные пространства над полем~$\KK$ (или свободные модули над кольцом~$R$).

\subsection{С высоты птичьего полёта}

Возьмём их декартово произведение $V_1 \times \dots \times V_m$ и отождестви в нём элементы, которые удовлетворяют соотношениям полилинейности:
\[
	\begin{aligned}
		(v_1, \dots, v_i' + v_i'', \dots, v_m)
		&\sim
		(v_1, \dots, v_i', \dots, v_m)
		+
		(v_1, \dots, v_i'', \dots, v_m)
		;
		\\
		(v_1, \dots, \lambda v_i, \dots, v_m)
		&\sim
		\lambda
		\cdot
		(v_1, \dots, v_i, \dots, v_m)
	\end{aligned}
\]
для любого $\lambda \in R$ ($\lambda \in \KK$) (соотношения линейности для каждого из множителей при фиксированных остальных).

Тензорное произведение --- это пространство, получающееся после такого отождествления:
\[
	V_1 \otimes \dots \otimes V_m
	:=
	(V_1 \times \dots \times V_m)/\sim
	,
\]
его элементы --- классы эквивалентностей $m$-ок векторов:
\[
	v_1 \otimes \cdots \otimes v_m := [(v_1, \dots, v_m)]_{\sim}.
\]

\subsection{Подпробнее}

Рассмотрим этот процесс детально.

Обозначим $V := V_1 \times \cdots \times V_m$.

Для начала, чтобы факторизовать декартово произведение по подпространству, нужно, чтобы в декартовом произведении была структура векторного пространства (или модуля над кольцом).
Ввести такую структуру означает научиться складывать элементы и умножать их на числа (точнее, элементы кольца). Сделаем сложение и умножение покомпонентным:
\begin{align*}
	(v_1', \dots, v_m') +
	(v_1'', \dots, v_m'') &:=
	(v_1' + v_1'', \dots, v_m' + v_m'')
	;
	\\
	k \cdot (v_1, \dots, v_m) &:=
	(kv_1, \dots, kv_m)
	.
\end{align*}
При этом нейтральным элементом по сложению будет набор вида~$(0, \dots, 0)$.

Если читательница сего опуса в достаточной степени любопытна (на что мы втайне надеемся), то может возникнуть вопрос: а почему определение \emph{именно такое}? Ничто же не мешает определить сложение как-нибудь поэкзотичнее --- к примеру, так:
\[
	(v_1', \dots, v_{m-1}', v_m') +
	(v_1'', \dots, v_{m-1}'', v_m'') :=
	(v_2' + v_2'', \dots, v_m' + v_m'', v_1' + v_1'')
	.
\]
Почему же был выбран другой вариант?

Один из способов придать смысл первоначальному выбору --- это такая коммутативная диаграмма:
\[
	\xymatrix{
		\left(
			\prod_{\mu = 1}^m V_\mu
		\right)^2
		\ar @{-->} [r] ^-{+}
		\ar [d] _-{\pi_\nu \times \pi_\nu}
		&
		\prod_{\mu = 1}^m V_\mu
		\ar [d] _-{\pi_\nu}
		\\
		V_\nu^2
		\ar [r] _-{+}
		&
		V_\nu
	}
\]
Она говорит: хочется, чтобы операция сложения, показанная пунктирной стрелкой сверху, замыкала диаграмму, то есть действовала таким образом, чтобы $\nu$-ая компонента~$w_\nu$ результата сложения двух наборов была такой же, как результат сложения $\nu$-ых компонент исходных наборов:
\[
	\xymatrix{
		\bigl(
			(\dots, v_\mu', \dots),
			(\dots, v_\mu'', \dots)
		\bigr)
		\ar @{|-->} [r] ^-{+}
		\ar @{|->} [d] _-{\pi_\nu \times \pi_\nu}
		&
		(\dots, w_\mu, \dots)
		\ar @{|->} [d] _-{\pi_\nu}
		\\
		\bigl(
			v_\nu',
			v_\nu''
		\bigr)
		\ar @{|->} [r] _-{+}
		&
		\left(
			\parbox{0.5in}{%
				\centering
				$w_\nu$\\
				$v_\nu' + v_\nu''$
			}%
		\right)
	}
\]
Если описать эту же диаграмму формулами, то в ней написано, что $(\pi_\nu \times \pi_\nu) \circ {+} = {+} \circ \pi_\nu$, то есть что взять сначала $\nu$-ые элементы пары, а затем сложить, --- то же самое, что сначала сложить наборы, а затем взять $\nu$-ый элемент результата:
\begin{align*}
	\bigl(
		+ \circ
		(\pi_\nu \times \pi_\nu)
	\bigr)
	\bigl(
		(\dots, v_\mu', \dots),
		(\dots, v_\mu'', \dots)
	\bigr)
	&=
	\bigl(
		\pi_\nu
		\circ
		+
	\bigr)
	\bigl(
		(\dots, v_\mu', \dots),
		(\dots, v_\mu'', \dots)
	\bigr)
	\\
	+
	\bigl(
		\pi_\nu(\dots, v_\mu', \dots),
		\pi_\nu(\dots, v_\mu'', \dots)
	\bigr)
	&=
	\pi_\nu
	\bigl(
		(\dots, v_\mu', \dots)
		+
		(\dots, v_\mu'', \dots)
	\bigr)
	\\
	+
	\bigl(
		v_\nu',
		v_\nu''
	\bigr)
	&=
	\pi_\nu
	\bigl(
		\dots, w_\mu, \dots
	\bigr)
	\\
	\bigl(
		v_\nu'
		+
		v_\nu''
	\bigr)
	&=
	w_\nu
	.
\end{align*}
\emph{Voil\`{a}}, формула готова.

